% Generated draft: requires mathematical and source review.
\documentclass[11pt]{article}
\usepackage[margin=25mm]{geometry}
\usepackage{fontspec}
\setmainfont{Latin Modern Roman}
\usepackage{amsmath,amssymb,mathtools,unicode-math}
\setmathfont{Latin Modern Math}
\usepackage{enumitem,needspace,xurl,hyperref}
\setlength{\parindent}{0pt}
\setlength{\parskip}{5pt}
\setlength{\emergencystretch}{3em}
\newcommand{\N}{\mathbb{N}}
\newcommand{\Z}{\mathbb{Z}}
\newcommand{\Q}{\mathbb{Q}}
\newcommand{\R}{\mathbb{R}}
\newcommand{\C}{\mathbb{C}}
\newcommand{\F}{\mathbb{F}}
\newcommand{\A}{\mathbb{A}}
\newcommand{\PP}{\mathbb{P}}
\newcommand{\E}{\mathbb{E}}
\newcommand{\eps}{\varepsilon}
\newcommand{\dd}{\,\mathrm{d}}
\begin{document}
\section*{UnsolvedMath ID: 30002297}
\textbf{Source problem code:} catalogue entry
\subsection{Definitions and mathematical statement}
Identify \(\C^2\) with time--frequency space \(\R^2\times\R^2\) by \(z=x+i\omega\). Define
\[
A_3=\begin{pmatrix}1&2/5\\0&\sqrt{21}/5\end{pmatrix},\qquad
\Lambda_3=A_3(\Z+i\Z)^2
=\{A_3m+iA_3n:m,n\in\Z^2\}.
\]
To make the operator conventions explicit, consider the customary normalized Gaussian and time--frequency shifts
\[
g(t)=\sqrt{2}\,e^{-\pi\lVert t\rVert^2},\qquad
(\pi(x,\omega)g)(t)=e^{2\pi i\omega\cdot t}g(t-x),
\qquad t,x,\omega\in\R^2.
\]
Here \(\lVert g\rVert_{L^2(\R^2)}=1\), and \(L^2(\R^2)\) is taken over \(\C\). Under these conventions, the associated Gabor system is
\[
G(g,\Lambda_3)=\{\pi(A_3m,A_3n)g:m,n\in\Z^2\}.
\]
The problem is to prove or disprove that there exist constants \(0<a\le b<\infty\) such that, for every \(f\in L^2(\R^2)\),
\[
a\lVert f\rVert_2^2
\le \sum_{m,n\in\Z^2}
\left|\left\langle f,\pi(A_3m,A_3n)g\right\rangle\right|^2
\le b\lVert f\rVert_2^2.
\]
The supplied record specifies a normalized Gaussian and an associated Gabor system but does not explicitly state the Gaussian exponent, modulation normalization, or complex-to-phase-space convention. The displayed choices give a precise conventional interpretation; their agreement with the report requires human verification. In particular, unit norm alone does not determine Gaussian width.
\subsection{Short English statement}
Determine whether the normalized Gaussian Gabor system on the specified lattice \(\Lambda_3\) is a frame for \(L^2(\R^2)\), meaning that its coefficients satisfy uniform positive lower and finite upper frame bounds. The explicit Gaussian and time--frequency conventions in this draft require confirmation against the cited report.
\subsection{Sources}
\begin{itemize}
\item Supplied local registry record: permanent UnsolvedMath ID \(30002297\), source problem code \texttt{catalogue entry}, in the Oberwolfach Reports Open Problems collection. The record gives proposed year 2013 and names no individual proposer.
\item \emph{Structured Function Systems and Applications} (2013), Oberwolfach Reports, pp.\ 579--655, DOI \texttt{10.4171/owr/2013/11}. Supplied source and literature reference: \url{https://doi.org/10.4171/owr/2013/11}.
\item Additional supplied literature link, labelled \emph{Structured Function Systems and Applications}: \url{https://ems.press/content/serial-article-files/46446}.
\item The supplied statement status is \texttt{corrected\_verified}. Its verification note states that the normalized Gaussian, generating matrix, complex lattice, and ambient Hilbert space omitted by extraction were restored from the official report. This is supplied provenance, not an independent verification by this draft.
\item Literature assessment dated exactly \texttt{2026-08-21}: the supplied curation labels the problem open, reports that this specific two-dimensional complex-lattice Gaussian system was conjectured to be a frame in the report, and states that no reliable proof or counterexample for this exact lattice was found. This assessment is supplied dated curation derived from the OWR material, not independently verified here; time-sensitive status and cited claims require independent checking.
\end{itemize}
\end{document}
